\documentclass[11pt]{article}
\usepackage[utf8]{inputenc}
\usepackage[T1]{fontenc}
\usepackage[english]{babel}
\usepackage{booktabs}
\usepackage{hyperref}
\usepackage{amsmath}
\usepackage{geometry}
\usepackage{graphicx}
\usepackage{caption}
\usepackage{float}
\geometry{margin=1in}
\usepackage{tikz}
\usepackage{pgfplots}
\pgfplotsset{compat=1.18}

\title{
    \textbf{Probing the Universality of Speech SSL Representations:}\\
    \large From Low-Resource ASR to Paralinguistic Robustness \\
    \vspace{0.5cm}
    \small \textit{A multi-faceted evaluation following the ML-SUPERB protocol}
}

\author{
    \begin{tabular}{c @{\hspace{1.5cm}} c @{\hspace{1.5cm}} c}
        \textbf{Janis Aiad} & \textbf{Vadim Lagresle} & \textbf{Bruny Soutarson}
    \end{tabular} \\
    \vspace{0.4cm}
    \small Project Supervisor: Maxime Poli (ENS-PSL)
}
\date{}
\begin{document}

\maketitle 

% ==============================================================================
% PART I: JANIS
% ==============================================================================
\section*{Part I: Architectural Innovation -- Benchmarking Joint Embedding Predictive Architectures (JEPA)}
\textit{Contributor: Janis}

\subsection*{Objective: Beyond Masked Prediction}
The baseline of this project follows the \textbf{ML-SUPERB} protocol: freezing a pretrained SSL speech encoder and training a lightweight CTC ASR head for low-resource budgets (10 minutes and 1 hour of supervised data per language). The core innovation involves comparing \textbf{JEPA-style} frontends (Audio-JEPA and WavJEPA) to \textbf{HuBERT} on the same recipe. We explore whether the JEPA objective, which predicts latent states rather than raw waveforms or discrete IDs, provides a more robust representation for low-resource transfer.

\subsection*{Monolingual Benchmark Results}
The performance was first validated on the English (\texttt{eng1}) 10-minute slice. All values below represent test CER and WER from the ESPnet results.

\begin{table}[h]
    \centering
    \caption{Monolingual performance comparison on English (\texttt{eng1}, 10 min)}
    \begin{tabular}{@{}lcc@{}}
        \toprule
        \textbf{Model Representation} & \textbf{CER (\%)} & \textbf{WER (\%)} \\ \midrule
        HuBERT (frozen, S3PRL) & 33.33 & 24.14 \\
        JEPA (Minimal pre-training) & 62.22 & 44.83 \\
        WavJEPA (Hugging Face pretrained) & 33.33 & 24.14 \\
        \textbf{WavJEPA (Local pretrain checkpoint)} & \textbf{33.33} & \textbf{24.14} \\ \bottomrule
    \end{tabular}
\end{table}

The results show that HuBERT and WavJEPA achieve parity on this slice. This equivalence suggests that once a pre-training threshold is met, both masked-prediction and joint-embedding architectures provide high-quality "universal" features for downstream ASR. The significantly higher error rates of the \textbf{JEPA minimal} model demonstrate that, unlike HuBERT, the JEPA latent space requires substantial pre-training data to converge on linguistically useful features.

\subsection*{Multilingual Pipeline Results}
Due to partial corpus availability (focusing on \texttt{eng}, \texttt{deu}, \texttt{fra}), the multilingual experiments benchmarked HuBERT in pooled ASR-only and joint ASR+LID settings.

\begin{table}[h]
    \centering
    \caption{Multilingual ASR metrics (HuBERT ASR-only)}
    \begin{tabular}{@{}lcc@{}}
        \toprule
        \textbf{Setting} & \textbf{CER (\%)} & \textbf{WER (\%)} \\ \midrule
        Multilingual ASR, 10 min & 24.96 & 23.48 \\
        Multilingual ASR, 1 h & 20.76 & 18.30 \\ \bottomrule
    \end{tabular}
\end{table}

We observe a clear \textbf{scaling effect}: moving from 10 minutes to 1 hour of data results in a \textbf{16.8\% relative reduction in CER}. This highlights that even with a frozen frontend, the lightweight CTC head benefits significantly from the increased phonetic diversity and regularization provided by the multilingual pool. 

Regarding \textbf{Language Identification (LID)}, the 10-minute track is completed. Note that for "LID-only" runs, traditional WER/CER metrics are not reported as the reference word count is zero; performance is instead validated via specific classification scores in the logs. The \textbf{ASR+LID (10 min)} joint track yielded a CER of 26.33\%, confirming the feasibility of simultaneous transcription and language classification.

\subsection*{Progress and Next Steps}
The current advancement is estimated at \textbf{70\%}. While the monolingual benchmarks and 10-minute joint tracks are stable, the following steps are prioritized:
\begin{itemize}
    \item Finalize the \textbf{ASR+LID 1 h} and \textbf{LoRA} (PEFT) rows. The LoRA experiments are particularly critical to test if targeted adaptation of the encoder's projection layers allows JEPA to resolve language-specific phonetic nuances more effectively than the purely frozen setup.
    \item Complete the numerical \textbf{ABX evaluation} using the developed tooling to quantify the phonetic discriminability of JEPA vs. HuBERT representations independently of the ASR head.
    \item Explore long-horizon WavJEPA pre-training to determine the performance ceiling of the local checkpoint.
\end{itemize}
% ==============================================================================
% PART II: VADIM (USER)
% ==============================================================================

\newpage
\section*{Part II: Interpretability Analysis -- Probing Layer Activation Weights}
\textit{Contributor: Vadim}

\subsection{Objective and Scope}
As highlighted in the ML-SUPERB benchmark regarding the complexity of Asian languages and the instability of certain SSL models in very low-resource regimes, our initial objective was to include Mandarin (CMN) in this interpretability study to contrast with Western European language families. However, due to significant environment-related technical challenges and persistent bugs in the Mandarin data preparation pipeline, we shifted our focus to a comparative analysis between the findings of Shi et al.  on XLSR and our own results using HuBERT and JEPA.

\subsection{Methodology: Testing the XLSR Hypothesis}
The original ML-SUPERB paper performs a layer-wise weight analysis on the XLSR-128 model. Their study concludes that the most relevant layers for ASR are located in the middle of the encoder stack and explicitly states that the terminal layers are not the primary drivers of transcription performance. We conducted a similar analysis on our HuBERT and JEPA models to determine if this "middle-layer dominance" is a universal property of speech SSL or an architecture-specific trait.

\subsection{Results and Interpretation: The HuBERT-XLSR Divergence}
Our results reveal a stark contradiction to the behavior observed in XLSR models, suggesting that the underlying pre-training objective fundamentally reshapes the internal hierarchy of the latent space.

\begin{itemize}
    \item \textbf{Terminal Dominance in HuBERT:} Unlike the XLSR results, our HuBERT-base and Large models exhibit a "Terminal Dominance" effect, where Layer 24 receives the vast majority of the learnable weight. This implies that the downstream CTC head relies almost exclusively on the final representation. 
    \item \textbf{Architectural Interpretation:} This difference may be explained by the pre-training objectives. XLSR uses a contrastive learning approach. In contrast, HuBERT predicts discrete pseudo-labels generated via iterative clustering. This discrete prediction task likely forces HuBERT to consolidate complex lexical information only at the very end of the encoder stack, whereas XLSR's contrastive features are distributed more broadly.
    \item \textbf{Data Scaling (10 min vs 1 h):} We observe that increasing data to 1 hour (orange and brown curves in Fig. \ref{fig:layer_analysis}) triggers a slight activation jump in the mid-stack (layers 10--15). This suggests that with more supervision, the model begins to find "XLSR-like" phonetic cues in intermediate layers, though the final layer remains the primary bottleneck.
    \item \textbf{JEPA's Structural Uniformity:} The Audio-JEPA model shows a more uniform distribution of weights. This suggests that its predictive architecture, which focuses on latent representation prediction rather than discrete tokens, maintains a more balanced hierarchical structure. This makes JEPA a very interesting model in performance and in information spread accross layer for interpretation. A deepen analysis could been made only with JEPA and XLSR models to conduct our initial experience : compare layer activation among laguages. 
\end{itemize}

% ==============================================================================
% PART III: BRUNY
% ==============================================================================
\newpage
\section*{Part III: Paralinguistic Robustness and Neural Alignment}
\textit{Contributor: Bruny}

\subsection{Objective and Scope}
While ML-SUPERB extends the original SUPERB benchmark \cite{yang2021superb} to more diverse and low-resource settings, it does not explicitly include speech emotion recognition (SER). However, SER provides a complementary evaluation of SSL representations, as it probes paralinguistic information such as affect, which is not directly optimized by standard pre-training objectives. Studying SER therefore allows us to assess whether these models capture higher-level, non-linguistic cues beyond phonetic and lexical content.

\subsection{Methodology: Low-Resource SER and Cross-Dataset Evaluation}
We evaluate the base versions of three SSL models (wav2vec2, HuBERT, and WavJEPA) on SER using RAVDESS \cite{ravdess}. Following the SUPERB protocol, all backbones are frozen and only a shallow classification head is trained.

\subsection{Results and Interpretation}

Table~\ref{tab:ser_ravdess_results} in Appendix~\ref{sec:appendix} shows a consistent improvement when moving from the 10-minute to the full training regime across all models, confirming that additional supervision helps better exploit frozen SSL representations in a paralinguistic setting. The performance gap between HuBERT and wav2vec2 is consistent with observations from the SUPERB benchmark \cite{yang2021superb}, where HuBERT generally provides stronger representations under comparable conditions. WavJEPA achieves the best performance in both regimes, suggesting that its representation learning objective captures more informative features for emotion recognition. Overall, the results indicate a clear ranking (WavJEPA $>$ HuBERT $>$ wav2vec2).

To assess the robustness of these representations, we further evaluate cross-corpus generalization by training on RAVDESS and testing on IEMOCAP \cite{iemocap} (see Table~\ref{tab:ser_iemocap_results} in Appendix~\ref{sec:appendix}). All models exhibit a performance drop due to the domain shift between acted and spontaneous emotional speech. However, the relative ordering is preserved: WavJEPA remains the most robust, followed by HuBERT, while wav2vec2 shows the largest degradation. This indicates that representations learned with more structured or semantic objectives transfer better to out-of-domain emotional data.

Beyond aggregate SER performance, we probe the internal hierarchy of SSL representations by comparing early, middle, and late layers. For each backbone, we extract representations at three depths of the encoder while keeping the model frozen and training the same shallow classifier.

Results show that emotional information is not maximally captured in the earliest layers (see Table~\ref{tab:ser_layer_results} in Appendix~\ref{sec:appendix}). For all three models, early representations are consistently weaker than middle and late ones, indicating that SER depends on more contextualized acoustic-prosodic features rather than purely local signal properties. However, the contrast between middle and late layers remains small. In HuBERT and wav2vec2, middle and late representations perform similarly, while in WavJEPA the late layers remain at least as strong as the middle ones. Rather than showing a strict middle-layer peak, these results suggest that emotion-relevant information is concentrated in the upper half of the encoder, with WavJEPA preserving useful paralinguistic structure more uniformly across depth.
% ==============================================================================
% BIBLIOGRAPHY
% ==============================================================================
\newpage
\begin{thebibliography}{9}
\bibitem{shi2023mlsuperb} Shi, B., et al. (2023). ML-SUPERB: Multilingual Speech Universal PERformance Benchmark. \textit{Interspeech}.
\bibitem{tuncay2025audiojepa} Tuncay, L., et al. (2025). Audio-JEPA: Joint-Embedding Predictive Architecture for Audio Representation Learning. \textit{arXiv:2507.02915}.
\bibitem{luthra2025spidr} Luthra, M., Poli, M., LeCun, Y., et al. (2025). SpidR-Adapt: A Universal Speech Representation Model for Few-Shot Adaptation. \textit{arXiv:2512.21204}.
\bibitem{raugel2025neural} Raugel, A., King, J.-R., et al. (2025). Hierarchical and temporal alignment between neural activity and language models. \textit{arXiv:2512.01591}.
\bibitem{watanabe2018espnet} Watanabe, S., et al. (2018). ESPnet: End-to-End Speech Processing Toolkit. \textit{Interspeech}.
\bibitem{ravdess} Livingstone, S. R., \& Russo, F. A. (2018). The Ryerson Audio-Visual Database of Emotional Speech and Song (RAVDESS): A dynamic, multimodal set of facial and vocal expressions in North American English.

\bibitem{iemocap} Busso, C., Bulut, M., Lee, C.-C., Kazemzadeh, A., Mower, E., Kim, S., Chang, J. N., Lee, S., \& Narayanan, S. S. (2007). IEMOCAP: Interactive emotional dyadic motion capture database.

\bibitem{yang2021superb} Yang, S.-W., Chi, P.-H., Chuang, Y.-S., Lai, C.-I. J., Lakhotia, K., Lin, Y. Y., Liu, A. T., Shi, J., Chang, X., Lin, G.-T., et al. (2021). SUPERB: Speech processing Universal PERformance Benchmark.
\end{thebibliography}

\newpage
\section{Appendix}
\label{sec:appendix}



\subsection{Pipeline and Software Stack}
The technical implementation is built on the \textbf{ESPnet} framework (\texttt{ml\_superb/asr1}) using the \textbf{uv} package manager for dependency resolution. To maintain compatibility with modern CUDA 12 backends and \textbf{Python 3.13}, several manual patches were applied:
\begin{itemize}
    \item \textbf{S3PRL Patching:} Manual modification of the S3PRL library to bypass deprecated \texttt{torchaudio} dependencies.
    \item \textbf{WavJEPA Integration:} Development of a custom loader for local WavJEPA checkpoints (Lightning) within the ESPnet frontend.
    \item \textbf{Automation:} Implementation of orchestration scripts for sequential queuing (\texttt{run\_multi.sh}) and automated result collection (\texttt{collect\_asr\_results.py}).
\end{itemize}

\subsection*{Hardware and Runtime Performance}
All experiments were conducted on a single \textbf{NVIDIA L4 GPU} (\texttt{janis-gpuL4-48}). The total computational cost to date is approximately \textbf{80 hours of H100-equivalent compute}. 

\begin{table}[h]
    \centering
    \caption{Measured Wall-Clock Training Times (HuBERT Multilingual)}
    \begin{tabular}{@{}llcc@{}}
        \toprule
        \textbf{Task} & \textbf{Data} & \textbf{Wall Time (h)} & \textbf{Training Only (h)} \\ \midrule
        ASR-only & 10 min & $\approx$ 6.1 h & 5.9 h \\
        ASR-only & 1 h & $\approx$ 11.6 h & 11.2 h \\
        LID-only & 10 min & $\approx$ 5.1 h & 4.9 h \\
        LID-only & 1 h & $\approx$ 11.5 h & 11.2 h \\ \bottomrule
    \end{tabular}
\end{table}

The pipeline supports \textbf{checkpoint resuming} (\texttt{--resume true}), allowing the system to recover from interruptions without loss of progress.

\newpage
\subsection{Layerwise analysis}

\begin{figure}[h]
    \centering
    % Augmentation de la taille de la première image pour une meilleure lisibilité
    \includegraphics[width=0.8\textwidth]{resultat_final_mva.png}
    \caption{Importance of SSL Layers (Normalized Weights) for HuBERT and JEPA. Note the massive spike at the terminal layer for HuBERT variants.}
    \label{fig:layer_analysis}
\end{figure}

\begin{figure}[h!]
    \centering
    \includegraphics[width=0.6\textwidth]{article.png}
    \caption{Article original layerwise weight analysis of XLSR-128 model in the monolingual track.}
    \label{fig:heatmap}
\end{figure}

\subsection{Emotional results}
\begin{table}[H]
\centering
\caption{
Speech emotion recognition results on RAVDESS using frozen SSL backbones (base versions only) with a shallow classifier head. 
Results are reported as mean $\pm$ std over multiple seeds.
}
\label{tab:ser_ravdess_results}
\begin{tabular}{llcccc}
\toprule
\textbf{Model} & \textbf{Train} & \textbf{UAR $\uparrow$} & \textbf{Macro-F1 $\uparrow$} & \textbf{Acc $\uparrow$} & \textbf{Loss $\downarrow$} \\
\midrule
HuBERT   & 10 min & $0.6523 \pm 0.0221$ & $0.6210 \pm 0.0262$ & $0.6667 \pm 0.0139$ & $0.8171 \pm 0.0182$ \\
HuBERT   & Full   & $0.7305 \pm 0.0013$ & $0.7298 \pm 0.0008$ & $0.7396 \pm 0.0035$ & $0.6187 \pm 0.0015$ \\
WavJEPA  & 10 min & $0.6745 \pm 0.0084$ & $0.6740 \pm 0.0091$ & $0.6806 \pm 0.0078$ & $0.7302 \pm 0.0116$ \\
WavJEPA  & Full   & $\mathbf{0.7448 \pm 0.0062}$ & $\mathbf{0.7542 \pm 0.0057}$ & $\mathbf{0.7569 \pm 0.0068}$ & $\mathbf{0.5634 \pm 0.0094}$ \\
wav2vec2 & 10 min & $0.5247 \pm 0.0273$ & $0.4834 \pm 0.0324$ & $0.5278 \pm 0.0139$ & $1.0384 \pm 0.0299$ \\
wav2vec2 & Full   & $0.5768 \pm 0.0352$ & $0.5434 \pm 0.0433$ & $0.5764 \pm 0.0417$ & $0.9899 \pm 0.0554$ \\
\bottomrule
\end{tabular}
\end{table}

\begin{table}[H]
\centering
\caption{Cross-corpus SER results (full train on RAVDESS, test on IEMOCAP).}
\label{tab:ser_iemocap_results}
\begin{tabular}{lcccc}
\toprule
\textbf{Model} & \textbf{UAR $\uparrow$} & \textbf{Macro-F1 $\uparrow$} & \textbf{Acc $\uparrow$} & \textbf{Loss $\downarrow$} \\
\midrule
HuBERT   & $0.5521 \pm 0.0184$ & $0.5347 \pm 0.0198$ & $0.5663 \pm 0.0162$ & $0.9315 \pm 0.0226$ \\
WavJEPA  & $\mathbf{0.6013 \pm 0.0147}$ & $\mathbf{0.5892 \pm 0.0161}$ & $\mathbf{0.6148 \pm 0.0139}$ & $\mathbf{0.8724 \pm 0.0185}$ \\
wav2vec2 & $0.4218 \pm 0.0229$ & $0.3895 \pm 0.0254$ & $0.4382 \pm 0.0211$ & $1.1876 \pm 0.0317$ \\
\bottomrule
\end{tabular}
\end{table}

\begin{table}[H]
\centering
\caption{Layer-level SER analysis on RAVDESS (full training set). For each SSL backbone, we compare early, middle, and late representations using the same frozen setup and shallow classifier.}
\label{tab:ser_layer_results}
\begin{tabular}{llccc}
\toprule
\textbf{Model} & \textbf{Layer group} & \textbf{UAR $\uparrow$} & \textbf{Macro-F1 $\uparrow$} & \textbf{Acc $\uparrow$} \\
\midrule
HuBERT   & Early  & $0.695 \pm 0.010$ & $0.690 \pm 0.011$ & $0.702 \pm 0.009$ \\
HuBERT   & Middle & $\mathbf{0.729 \pm 0.006}$ & $\mathbf{0.731 \pm 0.006}$ & $\mathbf{0.738 \pm 0.007}$ \\
HuBERT   & Late   & $0.730 \pm 0.001$ & $0.730 \pm 0.001$ & $0.740 \pm 0.003$ \\
\midrule
WavJEPA  & Early  & $0.715 \pm 0.008$ & $0.720 \pm 0.009$ & $0.725 \pm 0.008$ \\
WavJEPA  & Middle & $\mathbf{0.742 \pm 0.005}$ & $\mathbf{0.750 \pm 0.005}$ & $\mathbf{0.753 \pm 0.006}$ \\
WavJEPA  & Late   & $0.745 \pm 0.006$ & $0.754 \pm 0.006$ & $0.757 \pm 0.006$ \\
\midrule
wav2vec2 & Early  & $0.555 \pm 0.024$ & $0.525 \pm 0.027$ & $0.562 \pm 0.022$ \\
wav2vec2 & Middle & $\mathbf{0.575 \pm 0.020}$ & $\mathbf{0.548 \pm 0.023}$ & $\mathbf{0.585 \pm 0.019}$ \\
wav2vec2 & Late   & $0.577 \pm 0.035$ & $0.543 \pm 0.043$ & $0.576 \pm 0.041$ \\
\bottomrule
\end{tabular}
\end{table}

\subsubsection{Metrics}
\label{app:metrics}

Let $\mathcal{D} = \{(x_i, y_i)\}_{i=1}^{N}$ denote an evaluation set of size $N$, where
$y_i \in \{0,1,2,3\}$ is the ground-truth emotion label for sample $i$ and
$\hat{y}_i \in \{0,1,2,3\}$ is the predicted label obtained by
\[
\hat{y}_i = \arg\max_{c \in \{0,1,2,3\}} z_{i,c},
\]
where $z_{i,c}$ is the logit produced by the classifier for class $c$.

In our setup, the classifier predicts one of four emotion classes:
\[
\mathcal{C} = \{\texttt{neutral}, \texttt{happy}, \texttt{sad}, \texttt{angry}\}.
\]

\paragraph{Accuracy.}
\[
\mathrm{Accuracy}
=
\frac{1}{N}
\sum_{i=1}^{N}
\mathbf{1}(\hat{y}_i = y_i).
\]

\paragraph{Unweighted Average Recall (UAR).}
Let
\[
\mathrm{Recall}_c
=
\frac{\mathrm{TP}_c}{\mathrm{TP}_c + \mathrm{FN}_c}.
\]
Then
\[
\mathrm{UAR}
=
\frac{1}{|\mathcal{C}|}
\sum_{c \in \mathcal{C}} \mathrm{Recall}_c.
\]

\paragraph{Macro-F1.}
Let
\[
\mathrm{Precision}_c
=
\frac{\mathrm{TP}_c}{\mathrm{TP}_c + \mathrm{FP}_c},
\qquad
\mathrm{F1}_c
=
\frac{2\,\mathrm{Precision}_c\,\mathrm{Recall}_c}
{\mathrm{Precision}_c + \mathrm{Recall}_c}.
\]
Then
\[
\mathrm{Macro\text{-}F1}
=
\frac{1}{|\mathcal{C}|}
\sum_{c \in \mathcal{C}} \mathrm{F1}_c.
\]

\paragraph{Cross-entropy loss.}
Let
\[
p_{i,c}
=
\frac{\exp(z_{i,c})}{\sum_{k \in \mathcal{C}} \exp(z_{i,k})}.
\]
The average loss over $\mathcal{D}$ is
\[
\mathcal{L}
=
-\frac{1}{N}
\sum_{i=1}^{N}
\log p_{i,y_i}.
\]

\end{document}